\documentclass[12pt,a4paper]{report}

% ============================================================================
% PACKAGES
% ============================================================================

\usepackage[T5]{fontenc}
\usepackage[utf8]{inputenc}
\usepackage[vietnamese]{babel}
\usepackage{lmodern}
\usepackage{microtype}
\usepackage{placeins}
\usepackage{geometry}
\geometry{
	top=3cm,
	bottom=3cm,
	left=3.5cm,
	right=2.5cm
}
\usepackage{setspace}
\onehalfspacing

\clubpenalty=10000
\widowpenalty=10000

\usepackage{amsmath}
\usepackage{amssymb}
\usepackage{amsthm}
\usepackage{mathtools}

\usepackage{booktabs}
\usepackage{adjustbox}
\usepackage{multirow}
\usepackage{array}

\usepackage{graphicx}
\usepackage{caption}
\usepackage{subcaption}
\captionsetup{font=small, labelfont=bf}

\usepackage[colorlinks=true, linkcolor=black, citecolor=black, urlcolor=blue]{hyperref}

\usepackage{indentfirst}
\usepackage{enumitem}
\usepackage{tocbibind}

\usepackage{tikz}
\usepackage{pgfplots}
\pgfplotsset{compat=1.18}
\usetikzlibrary{shapes.geometric, arrows.meta, positioning, fit,
	backgrounds, calc, decorations.pathreplacing}

% ============================================================================
% GLOBAL TIKZ DEFAULTS
% ============================================================================
\tikzset{
	every picture/.style={line width=0.8pt}
}

% ============================================================================
% THEOREM DEFINITIONS
% ============================================================================

\theoremstyle{plain}
\newtheorem{dinhnghia}{Định nghĩa}[chapter]
\newtheorem{dinhly}{Định lý}[chapter]
\newtheorem{baitoan}{Bài toán}[chapter]

\theoremstyle{definition}
\newtheorem{vidu}{Ví dụ}[chapter]

% ============================================================================
% MATHEMATICAL COMMANDS
% ============================================================================

\DeclareMathOperator*{\argmax}{argmax}
\DeclareMathOperator*{\argmin}{argmin}
\DeclareMathOperator*{\sigmoid}{sigmoid}

% ============================================================================
% DOCUMENT INFORMATION
% ============================================================================

\title{Báo Cáo Nghiên Cứu Kỹ Thuật AI: Bước Dịch Chuyển Từ Naive RAG
	Sang GraphRAG Dựa Trên Nguyên Lý Đồ Thị Của Neo4j}
\author{Ngô Đức Phan Tiến Đạt}
\date{2026}

% ============================================================================
% BEGIN DOCUMENT
% ============================================================================

\begin{document}
	
	\newcommand{\frontmatter}{\pagenumbering{roman}}
	\pagestyle{plain}
	
	% ============================================================================
	% FRONT MATTER — Title Page
	% ============================================================================
	
	\begin{titlepage}
		\centering
		\vspace*{1cm}
		
		\textbf{TRƯỜNG ĐẠI HỌC KHOA HỌC}\\
		\textbf{ĐẠI HỌC HUẾ}
		
		\vspace{2cm}
		
		\textbf{\Large KHOA CÔNG NGHỆ THÔNG TIN}
		
		\vspace{1.5cm}
		
		\includegraphics[width=2.5cm, keepaspectratio]{logo.png}\\[0.3cm]
		{\small\textit{Trường Đại Học Khoa Học -- Đại Học Huế}}
		
		\vspace{1.5cm}
		
		\textbf{\LARGE BÁO CÁO NGHIÊN CỨU KỸ THUẬT AI}
		
		\vspace{1cm}
		
		\textbf{\Large Bước Dịch Chuyển Từ Naive RAG Sang GraphRAG\\
			Dựa Trên Nguyên Lý Đồ Thị Của Neo4j}
		
		\vspace{2cm}
		
		\textbf{Sinh Viên:}\\[0.3cm]
		Ngô Đức Phan Tiến Đạt
		
		\vspace{1cm}
		
		\textbf{Năm:} 2026
		
		\vfill
	\end{titlepage}
	
	% ============================================================================
	% FRONT MATTER — Declaration
	% ============================================================================
	
	\thispagestyle{empty}
	\vspace*{3cm}
	
	\section*{LỜI CAM ĐOAN}
	
	\vspace{1cm}
	
	Em xin cam đoan rằng nội dung của báo cáo này là hoàn toàn nguyên bản và được
	thực hiện dưới sự hướng dẫn của giảng viên hướng dẫn trong lĩnh vực Trí Tuệ Nhân Tạo.
	
	Em xác nhận rằng:
	\begin{itemize}
		\item Không có phần nào của báo cáo này được sao chép từ các công trình đã công
		bố mà không có trích dẫn thích hợp.
		\item Tất cả các tham chiếu và trích dẫn đều được ghi rõ nguồn gốc.
		\item Các kết quả nghiên cứu được trình bày là trung thực và khách quan.
	\end{itemize}
	
	\vspace{2cm}
	
	\begin{flushright}
		Thành phố Huế, ngày \quad tháng \quad năm 2026\\[1cm]
		\textbf{Ngô Đức Phan Tiến Đạt}
	\end{flushright}
	
	\clearpage
	
	% ============================================================================
	% FRONT MATTER — Acknowledgment
	% ============================================================================
	
	\thispagestyle{empty}
	\vspace*{3cm}
	
	\section*{LỜI CẢM ƠN}
	
	\vspace{1cm}
	
	Em xin bày tỏ lòng biết ơn sâu sắc nhất đến các thầy giáo, cô Đoàn Thị Hồng
	Phước, khoa Công Nghệ Thông Tin, chuyên ngành Khoa Học Máy Tính, Trường Đại
	Học Khoa Học -- Đại Học Huế đã tận tình hướng dẫn và truyền đạt kiến thức
	trong suốt quá trình nghiên cứu.
	
	Đặc biệt, Em xin cảm ơn:
	\begin{itemize}
		\item Các nhà nghiên cứu hàng đầu trong cộng đồng GraphRAG và Neo4j đã mở ra
		hướng đi mới cho việc ứng dụng đồ thị tri thức trong hệ thống truy xuất
		và sinh văn bản.
		\item Các đồng nghiệp trong phòng thí nghiệm đã có những thảo luận giá trị
		và đóng góp ý kiến xây dựng.
		\item Gia đình và bạn bè đã luôn đồng hành và hỗ trợ trong suốt hành trình
		nghiên cứu.
	\end{itemize}
	
	Cuối cùng, Em xin cảm ơn tất cả những người đã gián tiếp và trực tiếp giúp đỡ
	Em hoàn thành báo cáo này.
	
	\vspace{2cm}
	
	\begin{flushright}
		Thành phố Huế, ngày \quad tháng \quad năm 2026
	\end{flushright}
	
	\clearpage
	
	% ============================================================================
	% FRONT MATTER — Abstract
	% ============================================================================
	
	\thispagestyle{empty}
	\vspace*{3cm}
	
	\section*{TÓM TẮT (Abstract)}
	\addcontentsline{toc}{section}{Tóm tắt (Abstract)}
	
	\vspace{1cm}
	
	Báo cáo kỹ thuật AI cấp độ tốt nghiệp này trình bày quá trình nghiên cứu hệ
	thống tư vấn tuyển sinh thông minh cho Trường Đại Học Khoa Học -- Đại Học Huế,
	sử dụng kiến trúc GraphRAG (Graph Retrieval-Augmented Generation) dựa trên
	sách đọc trên nền tảng neo4j là Essential GraphRAG của tác giả Tomaž Bratanič
	và các đồng nghiệp của mình được xuất bản.
	
	\textbf{Câu hỏi nghiên cứu:} Liệu việc thay thế kiến trúc Naive RAG truyền
	thống bằng GraphRAG có cải thiện đáng kể độ chính xác và khả năng suy luận
	đa bước trong bài toán hỏi đáp tuyển sinh đại học hay không?
	
	\textbf{Đóng góp chính của nghiên cứu:}
	\begin{enumerate}
		\item Phân tích toán học hệ thống về giới hạn của Naive RAG trong bài toán
		truy xuất thông tin đa bước (multi-hop retrieval).
		\item Thiết kế kiến trúc GraphRAG hoàn chỉnh cho dữ liệu tuyển sinh tiếng
		Việt với đồ thị tri thức Neo4j.
		\item Xây dựng bộ thực nghiệm so sánh định lượng giữa Naive RAG và GraphRAG
		trên bộ dữ liệu tuyển sinh HUSC.
	\end{enumerate}
	
	\textbf{Từ khóa:} GraphRAG, Neo4j, Knowledge Graph, Tư vấn tuyển sinh,
	Multi-hop Reasoning, Hybrid Retrieval.
	
	\clearpage
	
	% ============================================================================
	% FRONT MATTER — Table of Contents / Lists
	% ============================================================================
	
	\tableofcontents
	\clearpage
	
	\listoffigures
	\addcontentsline{toc}{section}{Danh mục hình vẽ}
	\clearpage
	
	\listoftables
	\addcontentsline{toc}{section}{Danh mục bảng biểu}
	\clearpage
	
	% ============================================================================
	% FRONT MATTER — Abbreviations
	% [FIX 5-minor] Xóa KNN khỏi bảng vì không xuất hiện trong nội dung 2 chương
	% ============================================================================
	
	\section*{DANH MỤC CÁC KÝ HIỆU VÀ CHỮ VIẾT TẮT}
	\addcontentsline{toc}{section}{Danh mục các ký hiệu và chữ viết tắt}
	
	\begin{table}[h]
		\centering
		\begin{tabular}{ll}
			\toprule
			\textbf{Ký hiệu / Viết tắt} & \textbf{Giải thích}                        \\
			\midrule
			RAG                         & Retrieval-Augmented Generation             \\
			GraphRAG                    & Graph-based Retrieval-Augmented Generation \\
			LLM                         & Large Language Model                       \\
			KG                          & Knowledge Graph (Đồ thị Tri thức)          \\
			RRF                         & Reciprocal Rank Fusion                     \\
			PPR                         & Personalized PageRank                      \\
			NER                         & Named Entity Recognition                   \\
			RE                          & Relation Extraction                        \\
			MRR                         & Mean Reciprocal Rank                       \\
			NDCG                        & Normalized Discounted Cumulative Gain      \\
			BM25                        & Best Matching 25                           \\
			NLP                         & Natural Language Processing                \\
			HUSC                        & Trường Đại Học Khoa Học -- Đại Học Huế    \\
			\bottomrule
		\end{tabular}
	\end{table}
	
	\clearpage
	
	% ============================================================================
	% MAIN BODY
	% ============================================================================
	
	\newcommand{\mainmatter}{\pagenumbering{arabic}}
	
	% ============================================================================
	% CHAPTER 1 — Naive RAG
	% ============================================================================
	
	\chapter{Cơ sở lý thuyết và Giới hạn của Naive RAG}
	\label{chap:naive_rag}
	
	\section{Động lực: Hạn chế Cốt lõi của Mô hình Ngôn ngữ Lớn}
	\label{sec:llm_limitations}
	
	Trước khi phân tích kỹ thuật của Naive RAG, cần thiết lập nền tảng động lực
	căn bản: tại sao các kiến trúc tăng cường truy xuất lại trở thành một hướng
	nghiên cứu tất yếu? Câu trả lời nằm ở bốn hạn chế cố hữu của chính LLM ---
	những hạn chế mà không thể khắc phục bằng cách mở rộng quy mô mô hình hay
	kéo dài thời gian huấn luyện~\cite{neo4j2024essential}.
	
	\textbf{Ngưỡng kiến thức (Knowledge Cutoff).} LLM được huấn luyện trên tập
	dữ liệu tĩnh được thu thập đến một thời điểm nhất định. Mọi thông tin phát
	sinh sau thời điểm đó --- chính sách tuyển sinh mới nhất, quy chế sửa đổi,
	chỉ tiêu cập nhật --- đều hoàn toàn nằm ngoài tầm nhận thức của mô hình, bất
	kể độ lớn hay kiến trúc của nó.
	
	\textbf{Thông tin lỗi thời.} Ngay trong vùng thời gian trước ngưỡng kiến thức,
	tri thức của mô hình vẫn có thể không còn phản ánh thực tế do sự thay đổi liên
	tục của môi trường: nhân sự biến động, quy định được điều chỉnh, số liệu thống
	kê được cập nhật theo từng năm học.
	
	\textbf{Ảo giác (Hallucination).} Khác với cơ sở dữ liệu truyền thống vốn tra
	cứu bản ghi thực, LLM mã hóa tri thức dưới dạng phân phối xác suất trên không
	gian ngôn ngữ. Hệ quả là mô hình có thể tự tin sản sinh ra thông tin hoàn toàn
	bịa đặt --- từ trích dẫn học thuật không tồn tại đến tên ngành học hay mã ngành
	sai lệch --- mà không có bất kỳ tín hiệu cảnh báo nào gửi đến người dùng.
	
	\textbf{Thiếu dữ liệu nội bộ.} Thông tin đặc thù của từng tổ chức --- tài liệu
	tuyển sinh, quy chế nội bộ, cơ sở dữ liệu sinh viên --- chưa bao giờ là một
	phần của tập huấn luyện công khai, khiến LLM hoàn toàn không có khả năng lý
	luận trực tiếp trên những nguồn dữ liệu này.
	
	Phương án fine-tuning --- huấn luyện lại mô hình trên tập dữ liệu đặc thù ---
	thoạt nhìn có vẻ là giải pháp tự nhiên nhưng lại không thực tế: chi phí tính
	toán lớn, chu kỳ triển khai dài, và không có gì đảm bảo mô hình sẽ thực sự ghi
	nhớ thông tin mới theo cách có thể kiểm soát và kiểm chứng. RAG ra đời như một
	giải pháp bổ sung tri thức \emph{tại thời điểm suy luận} (inference-time), tách
	biệt hoàn toàn vòng đời của mô hình và vòng đời của dữ liệu: khi tri thức thay
	đổi, ta chỉ cần cập nhật knowledge base, không cần chạm vào mô hình.
	
	Trong kiến trúc RAG cơ bản (Naive RAG), tri thức đầu vào được xử lý dưới dạng
	các đoạn văn bản độc lập (chunks) và được biểu diễn qua các vector nhúng
	(embeddings) trong một không gian liên tục. Nghiên cứu của Grossmann và cộng
	sự~\cite{docgraphrag2025} về Document GraphRAG cho thấy, giả sử $\mathcal{T}$
	là tập hợp các đoạn văn bản, ta có ánh xạ từ không gian ngữ nghĩa sang không
	gian vector $\mathbb{R}^d$ qua bộ mã hóa $\phi$:
	
	\begin{equation}
		\phi: \mathcal{T} \to \mathbb{R}^d
		\label{eq:embedding_map}
	\end{equation}
	
	Khi hệ thống tiếp nhận một truy vấn $q$, nó tiến hành tìm kiếm $k$ lân cận
	gần nhất (k-NN) dựa trên độ tương đồng Cosine (Cosine Similarity) giữa vector
	truy vấn $\mathbf{v}_q = \phi(q)$ và các vector tài liệu $\mathbf{v}_d = \phi(d)$
	theo công thức~\eqref{eq:cosine_similarity}:
	
	\begin{equation}
		S_C(\mathbf{v}_q, \mathbf{v}_d)
		= \frac{\mathbf{v}_q \cdot \mathbf{v}_d}
		{\|\mathbf{v}_q\|_2\,\|\mathbf{v}_d\|_2}
		= \cos(\theta)
		\label{eq:cosine_similarity}
	\end{equation}
	
	Nghiên cứu của Edge và cộng sự tại Microsoft~\cite{edge2024local} chỉ ra Naive
	RAG bộc lộ điểm yếu cốt lõi khi đối mặt với các truy vấn cần suy luận đa bước
	(multi-hop reasoning). Khi chia nhỏ tài liệu, cấu trúc ngữ nghĩa ban đầu bị
	phá vỡ, tạo ra hiện tượng phân mảnh ngữ cảnh (context fragmentation). Hệ thống
	lưu giữ từng đoạn thông tin rời rạc nhưng đánh mất sợi dây liên kết mạch lạc
	giữa chúng. Hình~\ref{fig:naive_rag_fragmentation} minh họa hiện tượng này.
	
	\begin{figure}[htbp]
		\centering
		\begin{adjustbox}{max width=\linewidth}
			\begin{tikzpicture}[
				node distance=1.2cm,
				box/.style={rectangle, rounded corners=3pt, draw=black!70,
					fill=blue!10, minimum width=2.2cm, minimum height=0.8cm,
					align=center, font=\small},
				arrow/.style={-Stealth, thick}
				]
				\node[box, fill=orange!20] (doc1) at (0,0)    {Doc A};
				\node[box, fill=orange!20] (doc2) at (0,-1.2) {Doc B};
				\node[box, fill=orange!20] (doc3) at (0,-2.4) {Doc C};
				\node[box, fill=orange!20] (doc4) at (0,-3.6) {Doc D};
				
				\node[box, fill=green!15]  (vec1) at (3.5,0)    {$\mathbf{v}_A$};
				\node[box, fill=green!15]  (vec2) at (3.5,-1.2) {$\mathbf{v}_B$};
				\node[box, fill=green!15]  (vec3) at (3.5,-2.4) {$\mathbf{v}_C$};
				\node[box, fill=green!15]  (vec4) at (3.5,-3.6) {$\mathbf{v}_D$};
				
				\draw[arrow] (doc1) -- node[above,font=\tiny]{$\phi$} (vec1);
				\draw[arrow] (doc2) -- (vec2);
				\draw[arrow] (doc3) -- (vec3);
				\draw[arrow] (doc4) -- (vec4);
				
				\node[box, fill=yellow!20] (query) at (6.5,-1.8) {$\mathbf{v}_q$};
				
				\draw[arrow, dashed, gray] (query) -- (vec1);
				\draw[arrow, dashed, gray] (query) -- (vec2);
				\draw[arrow, dashed, gray] (query) -- (vec3);
				
				\node[above, font=\footnotesize\bfseries] at (0,0.5)   {Tài liệu};
				\node[above, font=\footnotesize\bfseries] at (3.5,0.5) {Embedding};
				\node[above, font=\footnotesize\bfseries] at (6.5,0.5) {Truy vấn};
				
				\begin{scope}[on background layer]
					\node[draw=orange!60, dashed, rounded corners,
					fit=(doc1)(doc2)(doc3)(doc4),
					label={[font=\tiny,text=orange!80]above:Chunks độc lập}] {};
				\end{scope}
			\end{tikzpicture}
		\end{adjustbox}
		\caption{Sự đứt gãy ngữ nghĩa khi chia nhỏ tài liệu trong Naive RAG. Mỗi
			chunk được biểu diễn độc lập trong không gian vector; không có kết nối ngữ
			nghĩa nào được duy trì giữa chúng.}
		\label{fig:naive_rag_fragmentation}
	\end{figure}
	
	% -----------------------------------------------------------------------
	% [FIX 1 & 2] Sửa công thức isotropic limit (định nghĩa sigma^2 tường minh)
	% và thay thế công thức multi-hop decay sai bằng giới hạn góc đúng
	% -----------------------------------------------------------------------
	
	Hạn chế toán học sâu xa của Naive RAG xuất phát từ hai tính chất của không
	gian vector nhúng chiều cao. Tính chất thứ nhất là \emph{giới hạn đẳng hướng}
	(isotropic concentration). Giả sử các thành phần của mỗi vector nhúng
	$\mathbf{v}_i \in \mathbb{R}^d$ là độc lập và đồng phân phối (i.i.d.) với
	trung bình $\mu = 0$ và phương sai $\sigma^2 > 0$ trên mỗi chiều. Khi đó,
	theo Luật Số Lớn, khoảng cách bình phương kỳ vọng giữa hai vector bất kỳ hội
	tụ theo~\cite{ethayarajh2019contextual}:
	
	\begin{equation}
		\mathbb{E}\!\left[\|\mathbf{v}_i - \mathbf{v}_j\|^2\right]
		= 2d\sigma^2
		\;\longrightarrow\; \infty
		\quad\text{khi } d \to \infty
		\label{eq:isotropic_limit}
	\end{equation}
	
	Hệ quả là sau khi chuẩn hóa về vector đơn vị, mọi cặp vector có xu hướng
	phân bố đồng đều trên mặt cầu $S^{d-1}$, khiến phần lớn giá trị cosine tập
	trung gần $0$. Độ tương đồng cosine do đó dần mất khả năng phân biệt.
	
	Tính chất thứ hai --- nghiêm trọng hơn cho bài toán suy luận đa bước ---
	là \emph{sự thiếu bắc cầu} (non-transitivity) của cosine similarity.
	Đây là hạn chế cấu trúc: ngay cả khi mọi cặp thực thể liền kề trong chuỗi
	$A \to B \to C \to \cdots \to N$ có độ tương đồng cao, không có bảo đảm nào
	rằng $A$ và $N$ sẽ gần nhau trong không gian vector.
	
	Phân tích hình học cho thấy điều này rõ ràng. Với các vector đơn vị, góc
	$\theta_{ij} = \arccos\!\bigl(S_C(\mathbf{v}_i, \mathbf{v}_j)\bigr)$ chính là
	khoảng cách trắc địa (geodesic distance) trên mặt cầu đơn vị $S^{d-1}$ và
	thỏa mãn bất đẳng thức tam giác. Do đó, góc tích lũy sau $n$ bước nhảy bị
	chặn trên bởi tổng các góc thành phần:
	
	\begin{equation}
		\theta_{A,\,A+n}
		\;\leq\;
		\sum_{k=0}^{n-1} \theta_{A+k,\,A+k+1}
		\label{eq:angle_triangle}
	\end{equation}
	
	Đặt $\rho = \min_{k}\, S_C(\mathbf{v}_{A+k},\,\mathbf{v}_{A+k+1})$ là độ
	tương đồng bước nhảy nhỏ nhất, và $\theta_0 = \arccos(\rho)$ là góc lệch
	tương ứng. Từ~\eqref{eq:angle_triangle} và tính đơn điệu giảm của hàm
	$\cos(\cdot)$ trên $[0,\pi]$, ta suy ra giới hạn dưới sau:
	
	\begin{equation}
		S_C(\mathbf{v}_A,\,\mathbf{v}_{A+n})
		\;\geq\;
		\cos\!\bigl(n \cdot \arccos(\rho)\bigr)
		\label{eq:multi_hop_lower_bound}
	\end{equation}
	
	Công thức~\eqref{eq:multi_hop_lower_bound} cho thấy giới hạn dưới này suy
	giảm theo $n$ và trở về $0$ (rồi âm) ngay khi $n > \frac{\pi/2}{\arccos(\rho)}$.
	Điều này có nghĩa là kể cả trong trường hợp tốt nhất, hệ thống không còn có
	bảo đảm dương nào về sự tương đồng sau đủ nhiều bước nhảy. Với $\rho = 0{,}95$
	(độ tương đồng bước rất cao), ngưỡng này là $n^* = \lfloor\pi/(2\arccos(0{,}95))\rfloor
	\approx 5$ bước --- một con số rất nhỏ so với các chuỗi suy luận thực tế.
	Hình~\ref{fig:cosine_decay} minh họa sự suy giảm này.
	
	\begin{figure}[htbp]
		\centering
		\begin{tikzpicture}
			\begin{axis}[
				width=9cm, height=5.5cm,
				xlabel={Số bước nhảy ($n$)},
				ylabel={Giới hạn dưới của $S_C$},
				xlabel style={font=\small},
				ylabel style={font=\small},
				ymin=-0.15, ymax=1.05,
				xmin=0, xmax=5,
				legend style={at={(0.97,0.97)}, anchor=north east, font=\small},
				grid=major,
				grid style={draw=gray!20},
				mark size=3pt,
				ytick={-0.1, 0, 0.2, 0.4, 0.6, 0.8, 1.0},
				]
				% Dữ liệu tính từ công thức lý thuyết: cos(n * arccos(0.95))
				% n=0: 1.000; n=1: 0.950; n=2: 0.802; n=3: 0.573; n=4: 0.296; n=5: -0.021
				\addplot[blue!70, mark=*, thick, mark options={fill=blue!70}]
				coordinates {(0,1.000)(1,0.950)(2,0.802)(3,0.573)(4,0.296)(5,-0.021)};
				\addlegendentry{$\cos(n\cdot\arccos(\rho))$ ($\rho=0.95$)}
			\end{axis}
		\end{tikzpicture}
		\caption{Suy giảm giới hạn dưới của độ tương đồng Cosine theo số bước nhảy
			trong truy xuất đa bước, tính theo công thức lý thuyết
			$\cos(n\cdot\arccos(\rho))$ với $\rho = 0{,}95$
			(xem công thức~\eqref{eq:multi_hop_lower_bound}).
			Đây là minh họa định lượng của kịch bản tốt nhất: ngay cả với độ tương đồng
			bước rất cao ($\rho=0{,}95$), giới hạn dưới đã xuống dưới $0$ tại $n=5$,
			chứng tỏ không gian vector mất hoàn toàn khả năng định hướng sau đủ nhiều bước.}
		\label{fig:cosine_decay}
	\end{figure}
	
	Khi giới hạn dưới $\cos(n\cdot\arccos(\rho))$ trở thành âm, vector không gian
	không còn cung cấp bất kỳ tín hiệu định hướng có ý nghĩa nào. Điều này lý giải
	tại sao Naive RAG thường thất bại khi phải tổng hợp thông tin từ nhiều nguồn để
	giải quyết các câu hỏi phức hợp: sự phi-bắc-cầu của cosine similarity là một
	giới hạn cấu trúc không thể khắc phục bằng cách tinh chỉnh mô hình hay tăng
	chiều không gian nhúng.
	
	% -------------------------------------------------------------------
	\section{Chiến lược Phân đoạn Văn bản: Đánh đổi Giữa Độ chính xác và Ngữ cảnh}
	\label{sec:chunking_tradeoffs}
	
	Một nguồn hạn chế thực tiễn quan trọng của Naive RAG nằm ngay ở bước tiền xử
	lý đầu tiên: phân đoạn văn bản (text chunking). Hai mục tiêu mâu thuẫn định
	hình mọi chiến lược phân đoạn: \emph{độ chính xác của vector nhúng} và
	\emph{tính đầy đủ của ngữ cảnh}. Đoạn văn bản càng ngắn, vector nhúng càng
	tập trung phản ánh một ý tưởng cục bộ --- và do đó tìm kiếm độ tương đồng cho
	kết quả sắc nét hơn. Nhưng đoạn quá ngắn sẽ cắt đứt ngữ cảnh mà LLM cần để
	tổng hợp câu trả lời ở bước tạo sinh. Ngược lại, đoạn dài bảo tồn tính mạch
	lạc ngữ cảnh nhưng làm loãng tín hiệu embedding, dẫn đến truy xuất kém phân
	biệt hơn.
	
	% -----------------------------------------------------------------------
	% [FIX 3] Thay ký hiệu tỷ lệ thuận ∝ bằng mô tả đơn điệu chính xác hơn
	% [FIX 4] Đổi tên tham số trọng số thành λ để tránh xung đột với α của PPR
	% -----------------------------------------------------------------------
	
	Gọi $f(\ell)$ và $g(\ell)$ lần lượt là Retrieval Precision và Context
	Completeness theo kích thước chunk $\ell$ (tính theo số token). Hai hàm này
	thể hiện quan hệ đơn điệu nghịch chiều nhau:
	\begin{equation}
		f'(\ell) < 0
		\quad\text{và}\quad
		g'(\ell) > 0
		\qquad \forall\,\ell > 0
		\label{eq:chunking_monotone}
	\end{equation}
	Mối quan hệ đơn điệu này được xác nhận qua thực nghiệm nhưng \emph{không}
	tuyến tính: dạng hàm cụ thể phụ thuộc vào phân phối nội dung corpus. Hàm
	mục tiêu hợp nhất với tham số trọng số $\lambda \in (0,1)$ là:
	\begin{equation}
		\mathcal{L}(\ell)
		= \lambda \cdot f(\ell) + (1-\lambda) \cdot g(\ell)
		\label{eq:chunking_objective}
	\end{equation}
	Điểm tối ưu $\ell^*$ thỏa mãn điều kiện bậc nhất
	$\lambda f'(\ell^*) + (1-\lambda)g'(\ell^*) = 0$ và phụ thuộc vào cả phân
	phối corpus lẫn tham số $\lambda$ phản ánh ưu tiên của ứng dụng.
	
	Một chiến lược tinh vi hơn nhằm phá vỡ sự đánh đổi cứng nhắc trên là kỹ thuật
	\emph{Parent Document Retriever}~\cite{neo4j2024essential}: duy trì hai cấp biểu
	diễn song song cho cùng một tài liệu. \emph{Child chunks} có kích thước nhỏ
	được dùng để tìm kiếm chính xác, nhưng khi kết quả được xác định, hệ thống trả
	về \emph{parent document} --- đoạn cha lớn hơn bao hàm chunk con đó --- để cung
	cấp ngữ cảnh đầy đủ cho tầng tạo sinh. Kỹ thuật này thừa nhận một nguyên lý
	quan trọng: đơn vị tối ưu cho \emph{truy xuất} và đơn vị tối ưu cho \emph{sinh
		văn bản} không nhất thiết trùng nhau.
	
	Trong bối cảnh hệ thống baseline PaddedRAG được đánh giá tại
	Chương~\ref{chap:experiments}, chiến lược phân đoạn cố định
	$\texttt{chunk\_size}=350$, $\texttt{overlap}=70$ là một xấp xỉ thực dụng nhưng
	không giải quyết được mâu thuẫn nêu trên. Đây là một điểm yếu cấu trúc cần
	được khắc phục khi chuyển sang kiến trúc GraphRAG --- nơi thông tin không còn
	bị ràng buộc trong ranh giới phân đoạn tùy tiện mà được tổ chức theo quan hệ
	ngữ nghĩa tường minh trong đồ thị tri thức.
	
	\begin{figure}[htbp]
		\centering
		\begin{adjustbox}{max width=0.88\linewidth}
			\begin{tikzpicture}
				\begin{axis}[
					width=10cm, height=6.5cm,
					xlabel={Kích thước chunk $\ell$ (tokens)},
					ylabel={Chất lượng tương đối},
					xlabel style={font=\small},
					ylabel style={font=\small},
					xmin=0, xmax=10,
					ymin=0, ymax=1.1,
					xtick={1,3,5,7,9},
					xticklabels={Rất nhỏ, Nhỏ, Trung bình, Lớn, Rất lớn},
					xticklabel style={font=\footnotesize},
					yticklabel style={font=\footnotesize},
					legend style={at={(0.5,1.05)}, anchor=south, legend columns=2,
						font=\small},
					grid=major,
					grid style={draw=gray!20},
					]
					\addplot[blue!70, thick, smooth]
					coordinates {(1,1.0)(2,0.88)(3,0.72)(4,0.56)(5,0.42)
						(6,0.30)(7,0.20)(8,0.13)(9,0.08)};
					\addlegendentry{$f(\ell)$: Độ chính xác truy xuất (giảm đơn điệu)}
					
					\addplot[orange!80, thick, smooth, dashed]
					coordinates {(1,0.05)(2,0.14)(3,0.30)(4,0.48)(5,0.63)
						(6,0.75)(7,0.85)(8,0.92)(9,0.97)};
					\addlegendentry{$g(\ell)$: Tính đầy đủ ngữ cảnh (tăng đơn điệu)}
					
					\addplot[green!40, fill=green!15, opacity=0.5]
					coordinates {(4,0)(4,1.1)(6,1.1)(6,0)} \closedcycle;
					
					\addplot[mark=star, mark size=5pt, only marks,
					mark options={fill=yellow!80, draw=black}]
					coordinates {(5,0.42)};
					\node[font=\footnotesize, above right] at (axis cs:5,0.42)
					{$\ell^*$};
					
					\node[font=\tiny, text=green!60!black, rotate=90]
					at (axis cs:5,0.85) {Vùng khuyến nghị};
				\end{axis}
			\end{tikzpicture}
		\end{adjustbox}
		\caption{Đánh đổi giữa độ chính xác truy xuất $f(\ell)$ (đơn điệu giảm)
			và tính đầy đủ ngữ cảnh $g(\ell)$ (đơn điệu tăng) theo kích thước chunk
			$\ell$. Điểm tối ưu $\ell^*$ nằm tại vùng giao thoa (nền xanh lá) --- nơi
			cả hai mục tiêu đạt mức chấp nhận được đồng thời.
			Các đường cong là minh họa định tính; dạng hàm cụ thể phụ thuộc vào corpus.}
		\label{fig:chunking_tradeoff}
	\end{figure}
	
	% ============================================================================
	% CHAPTER 2 — GraphRAG
	% ============================================================================
	
	\chapter{Kiến trúc GraphRAG và Nền tảng Lý thuyết}
	\label{chap:graphrag_theory}
	
	% ============================================================================
	% [FIX 5] Gộp mục 2.1 (Tổng quan) và mục 2.3 (Các biến thể) thành một mục
	% duy nhất, loại bỏ phần giới thiệu lặp lại
	% ============================================================================
	
	\section{Các Biến thể Kiến trúc GraphRAG và Hệ sinh thái}
	\label{sec:graphrag_ecosystem}
	
	GraphRAG (Graph-based Retrieval-Augmented Generation) là một kiến trúc tiên
	tiến trong lĩnh vực xử lý ngôn ngữ tự nhiên (NLP), được thiết kế để khắc phục
	những hạn chế cơ bản của Naive RAG đã phân tích tại Chương~\ref{chap:naive_rag}.
	Trong khi Naive RAG biểu diễn tri thức dưới dạng các đoạn văn bản rời rạc và
	thực hiện truy xuất dựa trên độ tương đồng vector, GraphRAG tổ chức thông tin
	theo cấu trúc đồ thị tri thức (Knowledge Graph -- KG), duy trì mối liên kết ngữ
	nghĩa giữa các thực thể. Sự khác biệt cốt lõi nằm ở khả năng suy luận đa bước:
	đồ thị tri thức duy trì các cạnh có hướng giữa các thực thể, cho phép hệ thống
	duyệt tường minh chuỗi liên kết $A \to B \to C$ thay vì dựa vào cosine
	similarity --- vốn bị chứng minh là phi-bắc-cầu ở Chương~\ref{chap:naive_rag}.
	
	Hệ sinh thái GraphRAG đã phát triển mạnh mẽ với nhiều biến thể kiến trúc, mỗi
	loại tập trung vào những khía cạnh cụ thể của bài toán. Bảng~\ref{tab:graphrag_variants_comparison}
	tóm tắt bốn biến thể chính trước khi đi vào phân tích từng loại.
	
	\begin{table}[htbp]
		\centering
		\caption{So sánh các biến thể kiến trúc GraphRAG}
		\label{tab:graphrag_variants_comparison}
		\begin{adjustbox}{max width=\linewidth}
			\begin{tabular}{lcccc}
				\toprule
				\textbf{Biến thể}
				& \textbf{Cơ chế cốt lõi}
				& \textbf{Suy luận đa bước}
				& \textbf{Cập nhật thời gian thực}
				& \textbf{Độ phức tạp}         \\
				\midrule
				Microsoft GraphRAG & Cộng đồng (Leiden) & Có & Không & Cao         \\
				HippoRAG           & PPR sinh học       & Có & Không & Trung bình  \\
				LightRAG           & Truy xuất hai cấp  & Có & Có    & Trung bình  \\
				SubgraphRAG        & Bộ ba đơn giản     & Có & Có    & Thấp        \\
				\bottomrule
			\end{tabular}
		\end{adjustbox}
	\end{table}
	
	\subsection{Microsoft GraphRAG: Community-Based Approach}
	\label{subsec:microsoft_graphrag}
	
	Microsoft GraphRAG~\cite{edge2024local} sử dụng thuật toán phát hiện cộng đồng
	Leiden để phân nhóm các thực thể liên quan chặt chẽ trong đồ thị tri thức. Chất
	lượng phân hoạch cộng đồng được đo lường bởi độ đo modularity $Q$ theo
	công thức~\eqref{eq:modularity}:
	\begin{equation}
		Q = \frac{1}{2m}
		\sum_{ij}\!\left[
		A_{ij} - \frac{k_i\,k_j}{2m}
		\right]\delta(c_i,\,c_j)
		\label{eq:modularity}
	\end{equation}
	trong đó $A_{ij}$ là phần tử ma trận kề, $k_i$ là bậc của đỉnh $i$,
	$m = \frac{1}{2}\sum_{ij}A_{ij}$ là tổng số cạnh, và $\delta(c_i, c_j)$ là
	hàm delta Kronecker bằng $1$ khi $i$ và $j$ thuộc cùng một cộng đồng.
	
	Sau khi phát hiện cộng đồng, LLM được sử dụng để tạo ra các bản tóm tắt cộng
	đồng (community summaries). Kiến trúc này hỗ trợ hai chế độ truy vấn:
	\emph{Local query mode} dành cho câu hỏi tập trung vào một thực thể cụ thể,
	kết hợp truy xuất vector và duyệt đồ thị; và \emph{Global query mode} dành cho
	câu hỏi tổng hợp, sử dụng quy trình map-reduce trên tất cả bản tóm tắt cộng
	đồng. Hình~\ref{fig:microsoft_graphrag_pipeline} minh họa pipeline xử lý.
	
	\begin{figure}[htbp]
		\centering
		\begin{adjustbox}{max width=\linewidth}
			\begin{tikzpicture}[
				node distance=0.8cm and 1.6cm,
				box/.style={rectangle, rounded corners=3pt, draw=black!65,
					fill=blue!10, minimum height=0.85cm,
					minimum width=2.6cm, align=center, font=\small},
				ans/.style={rectangle, rounded corners=3pt, draw=green!70,
					fill=green!15, minimum height=0.85cm,
					minimum width=2cm, align=center, font=\small},
				arrow/.style={-Stealth, thick}
				]
				\node[box, fill=blue!20]  (lq)  at (0,0)    {Truy vấn\\cục bộ};
				\node[box]                (er)  at (3.2,0)   {Truy xuất\\thực thể};
				\node[box]                (gt)  at (6.4,0)   {Duyệt đồ thị\\(k-hop)};
				\node[box]                (lc)  at (9.6,0)   {Ngữ cảnh\\cục bộ};
				
				\node[box, fill=orange!15] (gq)  at (0,-2)  {Truy vấn\\toàn cục};
				\node[box]                 (cs)  at (3.2,-2) {Tóm tắt\\cộng đồng};
				\node[box]                 (mr)  at (6.4,-2) {Map-Reduce\\tổng hợp};
				\node[box]                 (gc)  at (9.6,-2) {Ngữ cảnh\\toàn cục};
				
				\node[box, fill=red!10]   (llm) at (12,-1)  {LLM\\Generator};
				\node[ans]                (out) at (14.4,-1) {Câu trả\\lời};
				
				\draw[arrow] (lq) -- (er);
				\draw[arrow] (er) -- (gt);
				\draw[arrow] (gt) -- (lc);
				\draw[arrow] (gq) -- (cs);
				\draw[arrow] (cs) -- (mr);
				\draw[arrow] (mr) -- (gc);
				\draw[arrow] (lc) -| (llm);
				\draw[arrow] (gc) -| (llm);
				\draw[arrow] (llm) -- (out);
			\end{tikzpicture}
		\end{adjustbox}
		\caption{Pipeline xử lý truy vấn của Microsoft GraphRAG. Hai luồng xử lý
			song song --- cục bộ (trên) và toàn cục (dưới) --- hội tụ tại LLM Generator để
			tổng hợp câu trả lời cuối cùng.}
		\label{fig:microsoft_graphrag_pipeline}
	\end{figure}
	
	\subsection{HippoRAG: Biologically-Inspired Retrieval}
	\label{subsec:hipporag}
	
	HippoRAG~\cite{gutierrez2024hipporag} đề xuất cách tiếp cận lấy cảm hứng từ
	mô hình trí nhớ vùng hải mã (hippocampus). Thay vì thực hiện truy xuất đa bước
	lặp đi lặp lại, HippoRAG nén toàn bộ quá trình thành một thao tác đơn lẻ bằng
	Personalized PageRank (PPR). Công thức cập nhật PPR được định nghĩa theo
	\eqref{eq:hipporag_ppr}:
	
	% -----------------------------------------------------------------------
	% [FIX 4] Ký hiệu α (damping factor) trong PPR — tách biệt khỏi λ (chunking)
	% [FIX minor] Giải thích lý do chọn chuẩn hóa row-stochastic D^{-1}A
	% -----------------------------------------------------------------------
	\begin{equation}
		\mathbf{p}^{(t+1)}
		= \alpha\,\hat{\mathbf{A}}\,\mathbf{p}^{(t)}
		+ (1-\alpha)\,\mathbf{s}
		\label{eq:hipporag_ppr}
	\end{equation}
	trong đó $\hat{\mathbf{A}} = \mathbf{D}^{-1}\mathbf{A}$ là ma trận chuyển tiếp
	(transition matrix) được chuẩn hóa theo hàng (row-stochastic), đảm bảo tổng mỗi
	hàng bằng $1$ và giữ cho $\mathbf{p}^{(t)}$ là một phân phối xác suất hợp lệ
	qua các bước lặp. Dạng chuẩn hóa này được chọn thay vì dạng đối xứng
	$\mathbf{D}^{-1/2}\mathbf{A}\mathbf{D}^{-1/2}$ vì PPR mô phỏng quá trình
	\emph{random walk} có hướng: xác suất chuyển tiếp từ $i$ đến $j$ chỉ phụ
	thuộc vào bậc của $i$. Tham số $\alpha \in (0,1)$ là hệ số giảm dần
	(damping factor, thường $\alpha = 0{,}85$), và $\mathbf{s}$ là vector hạt
	nhân (seed vector) được khởi tạo từ các thực thể trích xuất của truy vấn $q$.
	
	\subsection{LightRAG: Real-time Graph Updates}
	\label{subsec:lightrag}
	
	LightRAG~\cite{guo2024lightrag} giải quyết điểm yếu quan trọng của các hệ
	thống GraphRAG thế hệ đầu: khả năng thích ứng với dữ liệu mới. Việc xây dựng
	lại toàn bộ đồ thị mỗi khi có tài liệu mới có chi phí $\mathcal{O}(|D|)$ tỷ
	lệ với tổng kích thước tập tài liệu, không thực tế với các hệ thống cập nhật
	liên tục.
	
	LightRAG giới thiệu cơ chế cập nhật tăng dần (incremental update): khi tài
	liệu mới được thêm vào, hệ thống chỉ trích xuất các thực thể và quan hệ mới,
	sau đó cập nhật đồ thị bằng cách chèn, xóa hoặc sửa đổi các node và cạnh
	tương ứng, giảm độ phức tạp từ $\mathcal{O}(|E|\log|V|)$ xuống
	$\mathcal{O}(|\Delta E|\log|V|)$ với $|\Delta E|$ là số cạnh thay đổi. Đồng
	thời, LightRAG sử dụng chiến lược truy xuất hai cấp (dual-level retrieval):
	ở cấp độ thấp (local), hệ thống truy xuất các thực thể và $k$-lân cận trực
	tiếp; ở cấp độ cao (global), hệ thống khám phá tri thức qua các mối quan hệ
	cộng đồng phức tạp.
	
	\subsection{SubgraphRAG: Simplicity as Strength}
	\label{subsec:subgraphrag}
	
	SubgraphRAG~\cite{subgraphrag2025} đưa ra quan điểm rằng sự đơn giản là chìa
	khóa cho hiệu quả. Thay vì triển khai các thành phần phức tạp như phát hiện
	cộng đồng hay bản tóm tắt cộng đồng, SubgraphRAG chỉ tập trung vào việc truy
	xuất các đồ thị con (subgraphs) liên quan trực tiếp đến truy vấn thông qua các
	bộ ba chủ thể-quan hệ-đối tượng (subject-relation-object triplets). Phân tích
	của Li và cộng sự cho thấy, đối với nhiều tác vụ, sự phức tạp bổ sung từ các
	thành phần cộng đồng không mang lại lợi ích đáng kể so với chi phí triển khai.
	SubgraphRAG vì vậy là lựa chọn hấp dẫn cho các ứng dụng cần sự đơn giản, hiệu
	quả, và dễ bảo trì.
	
	\subsection{Kiến trúc Agentic RAG: Linh hoạt qua Điều phối Công cụ}
	\label{subsec:agentic_rag}
	
	Các biến thể GraphRAG ở trên đều hoạt động theo một pipeline truy xuất tương
	đối cố định: câu hỏi $\to$ truy xuất $\to$ tạo sinh. Kiến trúc Agentic
	RAG~\cite{neo4j2024essential} phá vỡ sự tuyến tính này bằng cách đặt LLM vào
	vai trò agent có khả năng quan sát, lập kế hoạch và tự quyết định chiến lược
	truy xuất tối ưu cho từng câu hỏi cụ thể.
	
	Ba thành phần cốt lõi bao gồm:
	
	\textbf{Retriever Router.} LLM phân tích đặc điểm câu hỏi và lựa chọn công
	cụ phù hợp: vector search cho câu hỏi ngữ nghĩa, full-text search cho tên
	riêng, Cypher query cho câu hỏi có cấu trúc số học, hoặc kết hợp nhiều công
	cụ cho câu hỏi phức hợp.
	
	\textbf{Answer Critic.} Sau khi tạo ra câu trả lời sơ bộ, một LLM thứ hai
	(hoặc cùng LLM trong một prompt khác) đánh giá chất lượng câu trả lời. Nếu
	chưa đủ căn cứ, vòng lặp truy xuất mới được kích hoạt.
	
	\textbf{Vòng lặp phản hồi.} Toàn bộ quy trình: Lập kế hoạch $\to$ Truy xuất
	$\to$ Tạo sinh $\to$ Đánh giá $\to$ (Lặp lại nếu cần). Số vòng lặp tối đa
	được giới hạn để kiểm soát chi phí tính toán.
	
	\begin{figure}[htbp]
		\centering
		\begin{adjustbox}{max width=0.88\linewidth}
			\begin{tikzpicture}[
				node distance=1.5cm,
				box/.style={
					rectangle, rounded corners=5pt, draw=black!60,
					minimum height=1.0cm, minimum width=3.4cm,
					align=center, font=\small, text width=3.2cm
				},
				arrow/.style={-Stealth, thick, draw=black!70},
				rarrow/.style={-Stealth, thick, draw=red!65, dashed},
				garrow/.style={-Stealth, thick, draw=green!55!black, line width=1.2pt}
				]
				\node[box, fill=purple!15, draw=purple!50]
				(query) at (0, 0)
				{Truy vấn\\người dùng $q$};
				
				\node[box, fill=green!20, draw=green!55!black, line width=1.2pt]
				(out) at (0, -7.5)
				{Câu trả lời\\cuối cùng};
				
				\node[box, fill=blue!12, draw=blue!50]
				(router) at (5, 0)
				{Retriever Router\\{\footnotesize LLM phân tích $q$}};
				
				\node[box, fill=orange!18, draw=orange!60]
				(critic) at (5, -5.0)
				{Answer Critic\\{\footnotesize Đánh giá chất lượng}};
				
				\node[box, fill=green!10, draw=green!50, text width=3.4cm]
				(tools) at (10, 0)
				{Tool Execution\\{\footnotesize Vector / BM25 / Cypher}};
				
				\node[box, fill=yellow!20, draw=yellow!60!black]
				(gen) at (10, -5.0)
				{LLM Generator\\{\footnotesize Tạo câu trả lời}};
				
				\draw[gray!50, dashed, rounded corners=8pt, line width=0.9pt]
				(3.2, 0.7) rectangle (11.8, -5.8);
				\node[font=\footnotesize\itshape, text=gray!70]
				at (7.5, -6.2) {Vòng lặp thích nghi};
				
				\draw[arrow] (query)  -- (router);
				\draw[arrow] (router) -- (tools);
				\draw[arrow] (tools)  -- (gen);
				\draw[arrow] (gen)    -- (critic);
				
				\draw[rarrow]
				(critic.west) -- ++(-0.8, 0)
				|- node[pos=0.25, left, font=\footnotesize, text=red!65]
				{Chưa đủ}
				node[pos=0.25, left, yshift=-0.4cm, font=\footnotesize, text=red!65]
				{-- thử lại}
				(router.west);
				
				\draw[garrow]
				(critic.south) -- ++(0, -0.8)
				-| node[pos=0.2, below, font=\footnotesize, text=green!55!black]
				{Đạt yêu cầu}
				(out.east);
			\end{tikzpicture}
		\end{adjustbox}
		\caption{Kiến trúc Agentic RAG với vòng lặp thích nghi. LLM đóng vai trò
			agent điều phối, tự động lựa chọn công cụ truy xuất phù hợp và đánh giá
			chất lượng câu trả lời; nếu chưa đạt ngưỡng, vòng lặp được kích hoạt lại
			với chiến lược điều chỉnh.}
		\label{fig:agentic_rag}
	\end{figure}
	
	% -------------------------------------------------------------------
	\section{Pipeline Xây dựng Đồ thị Tri thức}
	\label{sec:kg_pipeline}
	
	Quá trình xây dựng một đồ thị tri thức chất lượng cao là nền tảng then chốt
	cho hiệu quả của bất kỳ hệ thống GraphRAG nào. Pipeline này bao gồm ba giai
	đoạn chính: trích xuất thực thể và quan hệ, giải quyết thực thể (entity
	resolution), và thiết kế lược đồ (schema design) cho cơ sở dữ liệu đồ thị.
	
	\subsection{Trích xuất Thực thể và Quan hệ}
	\label{subsec:ner_relation_extraction}
	
	Giai đoạn đầu tiên trong pipeline là chuyển đổi dữ liệu văn bản phi cấu trúc
	thành các cấu trúc đồ thị có ý nghĩa thông qua hai nhiệm vụ chính: trích xuất
	thực thể đặt tên (Named Entity Recognition -- NER) và trích xuất quan hệ
	(Relation Extraction -- RE). Có ba phương pháp chính được sử dụng trong thực tế:
	
	Phương pháp dựa trên luật (rule-based NER) sử dụng các biểu thức chính quy
	và từ điển để nhận diện thực thể. Phương pháp này có độ chính xác cao đối với
	các mẫu cố định nhưng thiếu khả năng tổng quát hóa. Phương pháp tinh chỉnh mô
	hình (fine-tuned NER) sử dụng các mô hình ngôn ngữ như PhoBERT hoặc VinAI được
	huấn luyện đặc biệt trên tập dữ liệu tiếng Việt, đạt được sự cân bằng tốt giữa
	độ chính xác và khả năng tổng quát. Phương pháp dựa trên LLM (LLM-based
	extraction) sử dụng trực tiếp các mô hình ngôn ngữ lớn để trích xuất thực thể
	và quan hệ; mặc dù linh hoạt và mạnh mẽ, phương pháp này có thể gặp phải vấn
	đề ảo giác (hallucination) và chi phí tính toán cao.
	
	Hiệu quả của quá trình NER thường được đánh giá bằng chỉ số $F_1$:
	\begin{equation}
		F_1 = \frac{2 \times \text{Precision} \times \text{Recall}}
		{\text{Precision} + \text{Recall}}
		= \frac{2\,TP}{2\,TP + FP + FN}
		\label{eq:f1_ner}
	\end{equation}
	trong đó $TP$, $FP$, $FN$ lần lượt là số thực thể được nhận diện đúng (True
	Positive), nhận diện sai (False Positive), và bị bỏ sót (False Negative).
	
	\subsection{Entity Resolution và Coreference}
	\label{subsec:entity_resolution}
	
	Một thách thức lớn trong việc xây dựng đồ thị tri thức từ dữ liệu thực tế là
	vấn đề đồng nhất thực thể (entity resolution). Trong văn bản tuyển sinh của
	Trường Đại Học Khoa Học -- Đại Học Huế, cùng một thực thể có thể xuất hiện
	dưới nhiều hình thức: ``ĐHKH'', ``Trường ĐHKH'', và ``Đại Học Khoa Học'' đều
	chỉ cùng một trường đại học. Nếu không được giải quyết, vấn đề này sẽ dẫn đến
	việc tạo ra nhiều node riêng lẻ cho cùng một thực thể, làm suy yếu cấu trúc
	đồ thị và giảm hiệu quả truy xuất.
	
	Các phương pháp giải quyết bao gồm so khớp chuỗi (string matching) dựa trên
	khoảng cách chỉnh sửa (edit distance), so sánh độ tương đồng nhúng (embedding
	similarity), và liên kết dựa trên LLM. Ngưỡng tương đồng được áp dụng theo:
	\begin{equation}
		\text{merge}(e_i, e_j)
		= \mathbb{1}\!\left[
		S_C\!\left(\phi(e_i),\,\phi(e_j)\right) \geq \tau
		\right]
		\label{eq:entity_resolution}
	\end{equation}
	trong đó $\phi(\cdot)$ là hàm nhúng thực thể, $S_C$ là độ tương đồng Cosine
	theo~\eqref{eq:cosine_similarity}, và $\tau \in [0,1]$ là ngưỡng phân giải.
	Trong thực nghiệm của nghiên cứu này, giá trị $\tau = 0{,}92$ được chọn dựa
	trên bước kiểm thử (ablation study) trên tập xác thực của bộ dữ liệu tuyển
	sinh HUSC; giá trị cao phản ánh sự thận trọng cần thiết trước mật độ từ viết
	tắt lớn trong văn bản hành chính tiếng Việt.
	
	\subsection{Thiết kế Schema Neo4j cho Dữ liệu Tuyển sinh}
	\label{subsec:neo4j_schema}
	
	% -----------------------------------------------------------------------
	% [FIX minor] Cập nhật mô tả schema để nhất quán với hình vẽ:
	% phân biệt meta-relationships (HAS_ENTITY, IN_COMMUNITY) và
	% domain relationships (THUỘC_TRƯỜNG, CÓ_NGÀNH, ...) là các instance
	% của relationship type trong đồ thị
	% -----------------------------------------------------------------------
	
	Lược đồ (schema) của đồ thị tri thức Neo4j cho dữ liệu tuyển sinh được thiết
	kế dựa trên các khuyến nghị từ Essential GraphRAG~\cite{neo4j2024essential}.
	Các nhãn node (node labels) chính bao gồm \texttt{\_\_Entity\_\_},
	\texttt{\_\_Chunk\_\_}, và \texttt{\_\_Community\_\_}.
	
	Về kiểu quan hệ (relationship types), hệ thống phân biệt hai lớp rõ ràng:
	\begin{itemize}
		\item \textbf{Quan hệ meta-cấu trúc:} \texttt{HAS\_ENTITY} liên kết node
		\texttt{\_\_Chunk\_\_} với các thực thể nó đề cập; \texttt{IN\_COMMUNITY}
		gán thực thể vào cộng đồng tương ứng.
		\item \textbf{Quan hệ ngữ nghĩa miền:} là các cạnh nối giữa các thực thể,
		được đặt tên theo ngữ nghĩa của dữ liệu tuyển sinh, ví dụ:
		\texttt{THUỘC\_TRƯỜNG} nối Khoa vào Trường, \texttt{CÓ\_NGÀNH} nối Khoa
		với các ngành đào tạo. Các kiểu quan hệ này thuộc về lớp \texttt{RELATIONSHIP}
		trong schema tổng quát và được định nghĩa cụ thể trong quá trình xây dựng đồ thị.
	\end{itemize}
	
	Chiến lược lập chỉ mục (index strategy) bao gồm: vector index trên thuộc
	tính \texttt{embedding}, text index trên \texttt{name}, và composite index
	để tối ưu hóa các truy vấn phức tạp. Sơ đồ đầy đủ được minh họa trong
	Hình~\ref{fig:neo4j_schema}.
	
	\begin{figure}[htbp]
		\centering
		\begin{adjustbox}{max width=0.85\linewidth}
			\begin{tikzpicture}[
				node distance=1.6cm and 2.2cm,
				entity/.style={ellipse, draw=blue!70, fill=blue!15,
					minimum width=2.4cm, minimum height=0.9cm,
					align=center, font=\small},
				chunk/.style={rectangle, draw=green!70, fill=green!10,
					minimum width=2cm, minimum height=0.8cm,
					rounded corners=2pt, align=center, font=\small},
				community/.style={rectangle, draw=orange!70, fill=orange!10,
					rounded corners=6pt, minimum width=2.2cm,
					minimum height=0.8cm, align=center, font=\small},
				arrow/.style={-Stealth, thick, draw=black!60},
				rlabel/.style={font=\tiny, fill=white, inner sep=1pt}
				]
				\node[entity]    (univ)  at (0,0)     {Đại Học Huế};
				\node[entity]    (fac)   at (-2.5,-2) {Khoa CNTT};
				\node[entity]    (maj)   at (2.5,-2)  {Ngành CNTT};
				\node[chunk]     (ch1)   at (-2.5,-4) {Chunk 1};
				\node[chunk]     (ch2)   at (2.5,-4)  {Chunk 2};
				\node[community] (comm)  at (0,-3)    {Cộng đồng\\CNTT};
				
				% Quan hệ ngữ nghĩa miền (domain relationships)
				\draw[arrow] (fac) -- node[rlabel,above left]
				{THUỘC\_TRƯỜNG {\tiny[domain]}} (univ);
				\draw[arrow] (fac) -- node[rlabel,above]
				{CÓ\_NGÀNH {\tiny[domain]}} (maj);
				% Quan hệ meta-cấu trúc
				\draw[arrow] (ch1) -- node[rlabel,left]
				{HAS\_ENTITY {\tiny[meta]}} (fac);
				\draw[arrow] (ch2) -- node[rlabel,right]
				{HAS\_ENTITY {\tiny[meta]}} (maj);
				\draw[arrow] (fac) -- node[rlabel,left]
				{IN\_COMMUNITY {\tiny[meta]}} (comm);
				\draw[arrow] (maj) -- node[rlabel,right]
				{IN\_COMMUNITY {\tiny[meta]}} (comm);
			\end{tikzpicture}
		\end{adjustbox}
		\caption{Lược đồ Neo4j cho hệ thống tư vấn tuyển sinh HUSC. Các thực thể
			(ellipse xanh) được liên kết với chunk văn bản gốc (hình chữ nhật xanh lá)
			qua quan hệ meta-cấu trúc \texttt{HAS\_ENTITY}, được gán vào cộng đồng
			(hình chữ nhật cam) qua \texttt{IN\_COMMUNITY}, và kết nối với nhau qua
			các quan hệ ngữ nghĩa miền như \texttt{THUỘC\_TRƯỜNG}, \texttt{CÓ\_NGÀNH}.}
		\label{fig:neo4j_schema}
	\end{figure}
	
	\subsection{Truy vấn Có cấu trúc qua Text2Cypher}
	\label{subsec:text2cypher}
	
	Sau khi đồ thị tri thức được xây dựng và lưu trữ trong Neo4j, kỹ thuật
	Text2Cypher giải quyết bài toán giao diện người dùng bằng cách đặt LLM vào
	vai trò tầng dịch thuật trung gian --- chuyển đổi câu hỏi ngôn ngữ tự nhiên
	thành câu lệnh Cypher khả thi trên cơ sở dữ liệu đồ thị.
	
	Các thực hành tốt nhất trong triển khai Text2Cypher~\cite{neo4j2024essential}
	bao gồm bốn yếu tố then chốt: (1)~\emph{Few-shot examples} --- cung cấp từ 3
	đến 10 cặp ví dụ (câu hỏi tự nhiên, câu lệnh Cypher) trong prompt; (2)~\emph{Schema
		injection} --- đưa đầy đủ mô tả lược đồ đồ thị vào prompt; (3)~\emph{Terminology
		mapping} --- xây dựng bảng ánh xạ từ vựng người dùng sang tên kỹ thuật trong
	schema, ví dụ: ``ngành'' $\to$ nhãn node \texttt{Major}, ``chỉ tiêu'' $\to$
	thuộc tính \texttt{quota}; (4)~\emph{Output constraints} --- chỉ định định dạng
	đầu ra thuần Cypher để tạo điều kiện cho bước hậu xử lý tự động.
	
	Tầm quan trọng của Text2Cypher nằm ở khả năng xử lý các truy vấn yêu cầu lọc
	và so sánh trên thuộc tính có cấu trúc. Câu hỏi ``Ngành nào trong khoa Công
	nghệ Thông tin có điểm chuẩn năm 2024 cao hơn 24?'' đòi hỏi so sánh số học
	trên thuộc tính node --- điểm mạnh của Cypher nhưng nằm ngoài khả năng của
	cosine similarity.
	
	% -------------------------------------------------------------------
	\section{Chiến lược Truy xuất Lai (Hybrid Retrieval)}
	\label{sec:hybrid_retrieval}
	
	Một trong những điểm mạnh cốt lõi của các hệ thống GraphRAG hiện đại là khả
	năng kết hợp nhiều chiến lược truy xuất qua cơ chế truy xuất lai (Hybrid
	Retrieval), tận dụng ưu điểm của truy xuất từ khóa (BM25), truy xuất vector
	dày đặc (dense retrieval), và duyệt đồ thị (graph traversal).
	
	% -----------------------------------------------------------------------
	% [FIX minor] Thêm giả thiết "local tree-like structure" cho k-hop complexity
	% -----------------------------------------------------------------------
	
	Một kỹ thuật quan trọng là mở rộng lân cận $k$-bước ($k$-hop neighborhood
	expansion). Từ tập node hạt nhân (seed nodes), hệ thống duyệt qua tất cả các
	node có thể đến được trong vòng $k$ bước. Dưới giả thiết cấu trúc cây cục bộ
	(local tree-like structure, tức là ít chu trình ngắn), độ phức tạp của thao
	tác này được biểu diễn theo:
	\begin{equation}
		\mathcal{O}\!\left(|V_s| \cdot b^k\right)
		\label{eq:khop_complexity}
	\end{equation}
	trong đó $|V_s|$ là số node hạt nhân và $b$ là hệ số phân nhánh trung bình
	(average branching factor). Với đồ thị có chu trình, số node thực sự duyệt
	nhỏ hơn giới hạn này khi áp dụng đánh dấu ``đã thăm'' (visited marking).
	Đối với đồ thị tuyển sinh HUSC với $|V| \approx 2500$ và $b \approx 3$,
	truy xuất $k=3$ bước có độ phức tạp thô $\approx 67{,}500$ phép duyệt trước
	khi tỉa.
	
	Để kết hợp kết quả từ nhiều nguồn, thuật toán Reciprocal Rank Fusion
	(RRF)~\cite{cormack2009reciprocal} được sử dụng:
	\begin{equation}
		\text{RRF}(d) = \sum_{m \in M} \frac{1}{k_0 + r_m(d)}
		\label{eq:rrf}
	\end{equation}
	trong đó $r_m(d)$ là thứ hạng của tài liệu $d$ trong danh sách $m$ và $k_0$
	là hằng số điều chỉnh (thường $k_0 = 60$). Điểm tổng hợp có trọng số cuối
	cùng:
	\begin{equation}
		S_{\text{final}}(d,q)
		= w_{\text{BM25}}\,S_{\text{BM25}}
		+ w_{\text{Dense}}\,S_{\text{Dense}}
		+ w_{\text{Graph}}\,S_{\text{Graph}}
		\label{eq:weighted_fusion}
	\end{equation}
	trong đó $w_{\text{BM25}} + w_{\text{Dense}} + w_{\text{Graph}} = 1$.
	Bảng~\ref{tab:retrieval_strategies} so sánh hiệu quả tương đối của từng thành
	phần theo loại truy vấn.
	
	\begin{table}[htbp]
		\centering
		\caption{Hiệu quả tương đối của các thành phần truy xuất lai theo loại truy vấn}
		\label{tab:retrieval_strategies}
		\begin{tabular}{lccc}
			\toprule
			\textbf{Loại truy vấn} & \textbf{BM25} & \textbf{Dense} & \textbf{Graph} \\
			\midrule
			Từ khóa đơn giản      & Cao           & Trung bình     & Thấp           \\
			Ngữ nghĩa             & Thấp          & Cao            & Trung bình     \\
			Đa bước (multi-hop)   & Thấp          & Thấp           & Cao            \\
			\bottomrule
		\end{tabular}
	\end{table}
	
	\subsection{Truy vấn Lùi Bước (Step-back Prompting)}
	\label{subsec:stepback_prompting}
	
	Một điểm yếu tinh tế của truy xuất vector trực tiếp là hiện tượng \emph{độ
		đặc thù quá cao} (over-specificity): câu hỏi người dùng thường rất cụ thể,
	trong khi thông tin nền tảng để trả lời lại nằm ở mức độ trừu tượng cao hơn.
	
	Kỹ thuật Step-back Prompting~\cite{neo4j2024essential} giải quyết vấn đề này
	bằng cách bổ sung một bước tiền truy xuất: LLM tự động sinh ra câu hỏi tổng
	quát hơn $\hat{q}$ từ câu hỏi gốc $q$, sau đó thực hiện song song hai lần
	truy xuất:
	\begin{equation}
		C_{\text{final}} = \text{Retrieve}(q) \cup \text{Retrieve}(\hat{q})
		\label{eq:stepback_retrieval}
	\end{equation}
	trong đó $C_{\text{final}}$ là tập ngữ cảnh hợp nhất (sau loại trùng lặp)
	được cung cấp cho LLM sinh văn bản.
	
	Kỹ thuật này đặc biệt hiệu quả trong GraphRAG nhờ sự phân tách tự nhiên giữa
	Global Search (phù hợp với $\hat{q}$ tổng quát, dựa trên community summaries)
	và Local Search (phù hợp với $q$ cụ thể, dựa trên entity traversal). Trong
	ngữ cảnh tuyển sinh HUSC, ví dụ điển hình: câu hỏi gốc ``Ngành Công Nghệ
	Thông Tin có xét tuyển A00 không?'' được bổ sung bởi câu hỏi lùi bước ``Các
	phương thức xét tuyển áp dụng tại HUSC là gì?'' --- khai thác hai lớp tri
	thức bổ trợ nhau trong đồ thị.
	
		\begin{figure}[htbp]
		\centering
		\includegraphics[width=0.95\textwidth]{step_backprompt.png}
		\caption{Luồng xử lý của kỹ thuật Step-back Prompting.
			Câu hỏi gốc $q$ và câu hỏi lùi bước $\hat{q}$ được truy xuất
			song song qua hai chế độ bổ trợ nhau; tập ngữ cảnh hợp nhất
			$C_{\mathrm{final}}$ cho phép LLM tổng hợp câu trả lời với
			nền tảng tri thức đầy đủ hơn so với truy xuất đơn lẻ.}
		\label{fig:stepback_prompting}
	\end{figure}
	\FloatBarrier
	% -------------------------------------------------------------------
	\section{Vấn đề Kỹ thuật và Giải pháp}
	\label{sec:technical_challenges}
	
	Việc triển khai hệ thống GraphRAG trong thực tế đi kèm với bốn thách thức
	kỹ thuật chính.
	
	\textbf{Lỗi thời của đồ thị (graph staleness).} Khi cơ sở tri thức liên tục
	được cập nhật, đồ thị có thể nhanh chóng trở nên lỗi thời. LightRAG đề xuất
	cập nhật tăng dần; đối với các hệ thống không có cơ chế này, việc lên lịch
	xây dựng lại định kỳ là cần thiết nhưng tốn kém.
	
	\textbf{Ảo giác trong LLM (LLM hallucination).} LLM có thể tạo ra các thực
	thể hoặc quan hệ không tồn tại trong dữ liệu gốc khi thực hiện trích xuất.
	Để giảm thiểu, cần triển khai các cơ chế xác thực chéo (cross-validation)
	với nguồn dữ liệu đáng tin cậy.
	
	\textbf{Nút cổ chai về khả năng mở rộng (scalability bottleneck).} Trình lập
	kế hoạch truy vấn của Neo4j có thể gặp khó khăn với các đồ thị rất lớn. Việc
	tối ưu hóa chỉ mục vector chuyên dụng và viết các truy vấn Cypher hiệu quả là
	bắt buộc để đảm bảo hiệu suất.
	
	\textbf{Lỗi hợp nhất sai (false positive merges) trong entity resolution.}
	Khi ngưỡng $\tau$ trong~\eqref{eq:entity_resolution} được đặt quá thấp, hai
	thực thể khác nhau có thể bị coi là một, làm méo mó cấu trúc đồ thị. Giải
	pháp là kết hợp nhiều đặc trưng (tên, mô tả, ngữ cảnh) trong quá trình so
	sánh và sử dụng ngưỡng thận trọng.
	
	% -------------------------------------------------------------------
	% [FIX 7] Khung đánh giá vẫn đặt tại Chương 2 nhưng được đóng khung rõ ràng
	% là thiết kế đề xuất sẽ được kiểm chứng tại Chương 3
	% -------------------------------------------------------------------
	\section[Khung Đánh giá Đề xuất cho Hệ thống GraphRAG]
	{Khung Đánh giá Đề xuất cho Hệ thống \linebreak GraphRAG}
		
	Phần này trình bày \emph{thiết kế} khung đánh giá sẽ được áp dụng tại
	Chương 3. Mục đích việc này là xác lập các chỉ số và phương pháp
	đo lường trước khi tiến hành thực nghiệm, đảm bảo tính hệ thống và khả năng
	so sánh.
	
	Essential GraphRAG~\cite{neo4j2024essential} đề xuất ba chỉ số độc lập, mỗi
	chỉ số phản ánh một tầng khác nhau trong pipeline truy xuất--tạo sinh.
	
	\textbf{Context Recall (Độ bao phủ ngữ cảnh)} đo lường liệu bộ truy xuất có
	thực sự tìm được toàn bộ thông tin cần thiết để trả lời câu hỏi hay không.
	Chỉ số này thấp khi retriever bỏ sót các chunk, node hoặc subgraph quan trọng.
	
	\textbf{Faithfulness (Độ trung thực)} kiểm tra xem câu trả lời do LLM tạo ra
	có bám sát vào tập ngữ cảnh đã được truy xuất hay LLM đã bổ sung nội dung nằm
	ngoài nguồn cung cấp. Faithfulness có thể được ước lượng bằng độ tương đồng
	Cosine (theo~\eqref{eq:cosine_similarity}) giữa embedding của câu trả lời và
	embedding của tập ngữ cảnh; giá trị thấp chỉ ra ảo giác xảy ra ngay cả khi
	retrieval hoạt động tốt.
	
	\textbf{Answer Correctness (Độ chính xác câu trả lời)} đánh giá câu trả lời
	cuối cùng có đúng so với ground truth hay không --- chỉ số đầu cuối
	(end-to-end) quan trọng nhất từ góc độ người dùng.
	
	\begin{figure}[htbp]
		\centering
		\begin{tikzpicture}[
			font=\small,
			tier/.style={
				rectangle, rounded corners=6pt, draw,
				minimum width=11.0cm, minimum height=1.05cm,
				align=center, text width=10.8cm
			},
			hdr/.style={
				rectangle, rounded corners=4pt, draw=black!40,
				fill=gray!12, minimum width=3.2cm, minimum height=0.75cm,
				align=center, font=\footnotesize, text width=3.0cm
			},
			cell/.style={
				rectangle, rounded corners=4pt, draw,
				minimum width=3.2cm, minimum height=0.75cm,
				align=center, font=\footnotesize, text width=3.0cm
			},
			dcell/.style={
				rectangle, rounded corners=4pt, draw,
				minimum width=4.4cm, minimum height=0.75cm,
				align=center, font=\footnotesize, text width=4.2cm
			}
			]
			\node[tier, fill=blue!10, draw=blue!50!black] (t1) at (0, 0)
			{\textbf{Tầng 1 — Context Recall:}\quad
				Bộ truy xuất có tìm đủ thông tin cần thiết không?};
			
			\node[tier, fill=orange!13, draw=orange!55!black] (t2) at (0,-1.45)
			{\textbf{Tầng 2 — Faithfulness:}\quad
				Câu trả lời có bám sát ngữ cảnh đã truy xuất không?};
			
			\node[tier, fill=green!10, draw=green!50!black] (t3) at (0,-2.90)
			{\textbf{Tầng 3 — Answer Correctness:}\quad
				Câu trả lời cuối cùng có đúng không?};
			
			\draw[-Stealth, thick, gray!55]
			(-6.2,-3.35) -- (-6.2, 0.50)
			node[midway, left, font=\tiny\itshape, gray!65,
			align=center, text width=1.1cm]
			{Mức độ\\trừu tượng\\tăng dần};
			
			\node[font=\small\bfseries] at (0,-3.75) {Ma trận chẩn đoán (sẽ điền kết quả tại Chương 3};
			\draw[gray!40, dashed] (-5.8,-3.42) -- (5.8,-3.42);
			
			\node[hdr] at (-3.4,-4.60) {Context Recall};
			\node[hdr] at ( 0.3,-4.60) {Faithfulness};
			\node[dcell, fill=gray!12, draw=black!40]
			at ( 4.2,-4.60) {Chẩn đoán};
			
			\node[cell, fill=green!18, draw=green!55!black]
			at (-3.4,-5.55) {Cao \checkmark};
			\node[cell, fill=green!18, draw=green!55!black]
			at ( 0.3,-5.55) {Cao \checkmark};
			\node[dcell, fill=yellow!22, draw=yellow!65!black]
			at ( 4.2,-5.55) {Lỗi schema /\\entity resolution};
			
			\node[cell, fill=red!16, draw=red!50!black]
			at (-3.4,-6.50) {Thấp $\times$};
			\node[cell, fill=green!18, draw=green!55!black]
			at ( 0.3,-6.50) {Cao \checkmark};
			\node[dcell, fill=red!18, draw=red!50!black]
			at ( 4.2,-6.50) {LLM đang ảo giác\\(hallucination)};
			
			\node[cell, fill=green!18, draw=green!55!black]
			at (-3.4,-7.45) {Cao \checkmark};
			\node[cell, fill=red!16, draw=red!50!black]
			at ( 0.3,-7.45) {Thấp $\times$};
			\node[dcell, fill=orange!20, draw=orange!60!black]
			at ( 4.2,-7.45) {Lỗi tầng generation\\(temperature)};
		\end{tikzpicture}
		\caption{Khung đánh giá ba tầng đề xuất cho hệ thống GraphRAG. Mỗi chỉ số
			phản ánh một giai đoạn độc lập trong pipeline; ma trận chẩn đoán phía dưới
			cho phép xác định chính xác tầng nào gây ra suy giảm chất lượng tổng thể.
			Kết quả thực nghiệm cụ thể sẽ được trình bày tại Chương 3.}
		\label{fig:evaluation_3metrics}
	\end{figure}
	
	Quy trình đánh giá thực nghiệm cho bài toán HUSC, dựa trên khung ba chiều này,
	sẽ được triển khai tại Chương 3 trên bộ câu hỏi được cô Phước phân
	tầng theo loại truy vấn.
	
	\clearpage
	\addcontentsline{toc}{chapter}{Tài liệu tham khảo}
	
	\begin{thebibliography}{99}
		
		\bibitem{vaswani2017attention}
		Vaswani, A., Shazeer, N., Parmar, N., Uszkoreit, J., Jones, L., Gomez, A.~N.,
		\ldots{} \& Polosukhin, I. (2017).
		\textit{Attention is all you need}.
		In \textit{Advances in Neural Information Processing Systems} (NeurIPS), 30,
		5998--6008.
		
		\bibitem{lewis2020retrieval}
		Lewis, P., Perez, E., Piktus, A., Petroni, F., Karpukhin, V., Goyal, N.,
		\& Kiela, D. (2020).
		\textit{Retrieval-Augmented Generation for Knowledge-Intensive NLP Tasks}.
		In \textit{Advances in Neural Information Processing Systems} (NeurIPS), 33,
		9459--9474.
		
		\bibitem{edge2024local}
		Edge, D., Trivedi, H., Monga, V., Ordun, G., Bastian, M., Lally, J.,
		\& Korreck, J. (2024).
		\textit{From Local to Global: A GraphRAG Approach to Query-Focused Summarization}.
		arXiv preprint arXiv:2404.16130.
		
		\bibitem{blondel2008fast}
		Blondel, V.~D., Guillaume, J.~L., Lambiotte, R., \& Lefebvre, E. (2008).
		\textit{Fast unfolding of communities in large networks}.
		\textit{Journal of Statistical Mechanics: Theory and Experiment}, 2008(10), P10008.
		
		\bibitem{manning2008introduction}
		Manning, C.~D., Raghavan, P., \& Schütze, H. (2008).
		\textit{Introduction to Information Retrieval}.
		Cambridge University Press.
		
		\bibitem{neo4j2024essential}
		Neo4j Inc. (2024).
		\textit{Essential GraphRAG}. Technical Report.
		\url{https://go.neo4j.com/rs/710-RRC-335/images/Essential-GraphRAG.pdf}
		
		\bibitem{docgraphrag2025}
		Grossmann, D., Caymazer, O., \& Knollmeyer, S. (2025).
		\textit{Document GraphRAG: Knowledge Graph Enhanced Retrieval Augmented Generation
			for Document Question Answering Within the Manufacturing Domain}.
		\textit{MDPI Electronics}, 14(11), 2102.
		
		\bibitem{subgraphrag2025}
		Li, M., Miao, S., \& Li, P. (2025).
		\textit{Simple is Effective: The Roles of Graphs and Large Language Models in
			Knowledge-Graph-Based Retrieval-Augmented Generation}.
		In \textit{Proceedings of ICLR 2025}.
		
		\bibitem{grag2025}
		Hu, Y., et al. (2025).
		\textit{GRAG: Graph Retrieval-Augmented Generation}.
		In \textit{Findings of NAACL 2025}, 4145--4157.
		
		\bibitem{ethayarajh2019contextual}
		Ethayarajh, K. (2019).
		\textit{How Contextual are Contextualized Word Representations?}
		In \textit{Proceedings of EMNLP-IJCNLP 2019}, 55--65.
		
		\bibitem{robertson2009bm25}
		Robertson, S., \& Zaragoza, H. (2009).
		\textit{The Probabilistic Relevance Framework: BM25 and Beyond}.
		\textit{Foundations and Trends in Information Retrieval}, 3(4), 333--389.
		
		\bibitem{cormack2009reciprocal}
		Cormack, G.~V., Clarke, C.~L.~A., \& Buettcher, S. (2009).
		\textit{Reciprocal rank fusion outperforms Condorcet and individual rank learning methods}.
		In \textit{Proceedings of SIGIR 2009}, 758--759.
		
		\bibitem{chen2020simple}
		Chen, D. (2020).
		\textit{A Simple Yet Tough-to-Beat Baseline for Non-English Dense Retrieval}.
		arXiv preprint arXiv:2012.15516.
		
		\bibitem{gutierrez2024hipporag}
		Gutierrez, B.~J., et al. (2024).
		\textit{HippoRAG: Neurobiologically Inspired Long-Term Memory for Large Language Models}.
		arXiv preprint arXiv:2405.14831.
		
		\bibitem{guo2024lightrag}
		He, H., et al. (2024).
		\textit{LightRAG: Simple and Fast Retrieval-Augmented Generation}.
		arXiv preprint arXiv:2410.05779.
		
	\end{thebibliography}
	
\end{document}